\chapter{Study Scenario}
\label{app:scenario}

The scenario reproduced below was presented to all participants in German, regardless of interaction mode (Form, Ask, Chat). The English translation is provided for the international reader; the German original is the version actually shown.

\section{Original (German)}
\label{app:scenario-de}

\begin{verbatim}
Szenario: Mitarbeiterverwaltung (Arbeitszeit/Urlaub) & Parkplatzberechtigung

Sie sind Mitarbeitender namens A.I. und arbeiten für die APOR GmbH, ein
Softwareunternehmen mit ca. 50 Mitarbeitenden in Ulm. Sie handeln im Namen
des Unternehmens und sind im Szenario sowohl Verantwortlicher (Data
Controller) als auch Datenschutzbeauftragte:r (DPO/DSB).

Zur Organisation:

Rolle: Data Controller, DPO
Adresse: Mainstraße 67, Ulm
E-Mail: mail@apor.eu
Telefon: 012345678
Vertreter: (Sie), A. I., ai@apor.eu
DPO: Das sind auch Sie

Verarbeitung
Das Unternehmen verarbeitet personenbezogene Daten, um Mitarbeitende vom
Eintritt bis zum Austritt zu verwalten. Im Fokus stehen:

- Anlegen neuer Mitarbeitender
- Dokumentation von Arbeitszeiten
- Verwaltung von Urlaub und Abwesenheiten

Zusätzlich verwaltet APOR für berechtigte Mitarbeitende eine
Parkplatzberechtigung für einen Firmenparkplatz mit Schranke. Damit
berechtigte Mitarbeitende den Parkplatz nutzen können, wird ein geeignetes
Merkmal zur Zuordnung der Berechtigung hinterlegt (z. B. ein
Fahrzeugkennzeichen).

APOR arbeitet mit intern gehostetem Microsoft Office sowie internen Servern
(keine externe Cloud). Krankentage oder medizinische Angaben sind für diese
Verarbeitung nicht erforderlich.

Für die Verwaltung sind technische und organisatorische Maßnahmen
vorgesehen, wie z.B. Zugriffsbeschränkung nach Rollen (Need-to-know),
Passwortschutz (ggf. MFA), regelmäßige Backups, Zutrittskontrolle zu
Serverraum/Archiv und Vertraulichkeitsverpflichtung der berechtigten
Mitarbeitenden.

Alle in diesem Zusammenhang gespeicherten Informationen werden spätestens
3 Monate nach Austritt der jeweiligen Person gelöscht.
\end{verbatim}

\section{English Translation}
\label{app:scenario-en}

\begin{verbatim}
Scenario: Employee Administration (Working Time / Leave) & Parking
Authorisation

You are an employee named A.I., working for APOR GmbH, a software company
with approximately 50 employees in Ulm. You act on behalf of the company
and, in this scenario, are both the Data Controller and the Data
Protection Officer (DPO).

Organisation details:

Role: Data Controller, DPO
Address: Mainstrasse 67, Ulm
E-mail: mail@apor.eu
Phone: 012345678
Representative: (You), A. I., ai@apor.eu
DPO: Also you

Processing
The company processes personal data in order to administer employees from
onboarding to exit. The focus is on:

- Onboarding new employees
- Documentation of working hours
- Administration of leave and absences

In addition, APOR manages a parking authorisation for entitled employees
for a company car park with a barrier. To allow entitled employees to use
the car park, a suitable feature for assigning the authorisation is
recorded (e.g. a vehicle licence plate).

APOR uses internally hosted Microsoft Office and internal servers (no
external cloud). Sick days or medical information are not required for
this processing.

For administration, technical and organisational measures (TOMs) are
provided, such as role-based access restriction (need-to-know), password
protection (where appropriate, MFA), regular backups, physical access
control to the server room / archive, and a confidentiality obligation for
authorised employees.

All information stored in this context is deleted no later than 3 months
after the person's departure.
\end{verbatim}


\chapter{Survey Instruments}
\label{app:sus-rtlx-items}

This appendix reproduces the verbatim items of the two standardised survey instruments referenced from the Methodology chapter: the System Usability Scale (SUS) and the Raw NASA Task Load Index (R-TLX). The six-item Perceived AI Support instrument is bespoke to this study and is listed in full in the Methodology chapter; it is not reproduced here.

\section{System Usability Scale (SUS)}
\label{app:sus-items}

Five-point Likert scale, $1 = $ Strongly Disagree, $5 = $ Strongly Agree. Odd-numbered items are positively keyed; even-numbered items are negatively keyed. The canonical wording reproduced below is taken from Brooke~\cite{Brooke1996}.

\begin{enumerate}
    \item I think that I would like to use this system frequently.
    \item I found the system unnecessarily complex.
    \item I thought the system was easy to use.
    \item I think that I would need the support of a technical person to be able to use this system.
    \item I found the various functions in this system were well integrated.
    \item I thought there was too much inconsistency in this system.
    \item I would imagine that most people would learn to use this system very quickly.
    \item I found the system very cumbersome to use.
    \item I felt very confident using the system.
    \item I needed to learn a lot of things before I could get going with this system.
\end{enumerate}

\section{NASA Raw Task Load Index Subscales}
\label{app:rtlx-items}

Six subscales, each rated on the standard NASA TLX 21-point scale. The Raw TLX variant used in this study takes the unweighted mean across the six subscales rather than applying the pairwise-comparison weighting of the original NASA TLX. The subscale definitions reproduced below follow Hart~\cite{Hart2006}.

\begin{description}
    \item[Mental Demand] How mentally demanding was the task?
    \item[Physical Demand] How physically demanding was the task?
    \item[Temporal Demand] How hurried or rushed was the pace of the task?
    \item[Performance] How successful were you in accomplishing what you were asked to do?
    \item[Effort] How hard did you have to work to accomplish your level of performance?
    \item[Frustration] How insecure, discouraged, irritated, stressed, and annoyed were you?
\end{description}


\chapter{LLM-as-a-Judge Configuration}
\label{app:judge}

This appendix documents the inputs to the LLM-as-a-judge evaluation layer that scored the 113 retained ROPA documents. It contains the verbatim system and user prompts, the five-dimension rubric, and the model configuration needed to reproduce the evaluation.

\section{System and User Prompt}
\label{app:judge-prompt}

The system prompt below is sent verbatim as the \texttt{system} role on every evaluation call. It is taken from \texttt{python/judge/prompt.json} (key \texttt{judge.system}).

\begin{verbatim}
You are an expert GDPR compliance reviewer evaluating Records of
Processing Activities (ROPA) documents under Article 30 GDPR. The
documents are written by novices (non-experts, no prior GDPR training)
who described the same processing scenario using different LLM-assisted
tools. Your job is to judge each ROPA document against (a) the provided
scenario and (b) GDPR Article 30 requirements. Be strict but fair. Do not
reward verbosity. Do not penalise a document for merely restating
scenario facts. Respond with valid JSON only - no prose, no Markdown
fences, no commentary before or after the JSON object.
\end{verbatim}

The accompanying user-message template is sent on every call with \texttt{\{scenario\}} (the German original from Appendix~\ref{app:scenario-de}) and \texttt{\{document\}} (the participant ROPA document) substituted in:

\begin{verbatim}
You will evaluate one ROPA document against the scenario it is supposed
to cover.

=== SCENARIO (verbatim, German) ===
{scenario}
=== END SCENARIO ===

=== ROPA DOCUMENT TO EVALUATE ===
{document}
=== END DOCUMENT ===

Score the document on each of the following five metrics using a 1-5
Likert scale, and give a short rationale (max 2 sentences, English) for
each score. The scale for all metrics is:
1 = very poor, 2 = poor, 3 = adequate, 4 = good, 5 = excellent.

Metrics and what to assess:
- completeness: Coverage of Art. 30(1) GDPR required elements -
  controller identity and contact, DPO, purposes of processing, legal
  basis, categories of data subjects and personal data, categories of
  recipients, retention periods, and technical/organisational measures
  (TOMs). Missing elements lower the score.
- scenario_faithfulness: How accurately the document reflects the
  concrete facts of THIS scenario - APOR GmbH, Ulm (Mainstrasse 67), ~50
  employees, internally hosted Microsoft Office and internal servers (NO
  external cloud), no medical/sickness data, 3-month deletion after
  employee departure, parking lot authorisation using a distinguishing
  feature such as a licence plate. Generic ROPA boilerplate that ignores
  these specifics lowers the score.
- legal_correctness: Correct use of GDPR terminology and correct citation
  of legal basis (typically Art. 6(1)(b) performance of contract for
  employment data, Art. 6(1)(c) legal obligation where applicable, Art.
  6(1)(f) legitimate interest for parking authorisation). Wrong articles,
  invented legal bases, or confusing controller/processor/DPO roles lower
  the score.
- hallucination_inverse: Inverse score for invented facts not supported
  by the scenario (e.g. invented third-party processors, invented data
  categories like medical data, invented addresses, invented retention
  periods different from 3 months). 5 = no hallucinations, 1 = many
  invented facts.
- overall: Holistic judgement of how useful this document would be as a
  real Art. 30 ROPA record for this scenario, taking the four dimensions
  above into account.

Return a single JSON object with exactly this shape:
{
  "completeness":           {"score": <int 1-5>, "rationale": "<string>"},
  "scenario_faithfulness":  {"score": <int 1-5>, "rationale": "<string>"},
  "legal_correctness":      {"score": <int 1-5>, "rationale": "<string>"},
  "hallucination_inverse":  {"score": <int 1-5>, "rationale": "<string>"},
  "overall":                {"score": <int 1-5>, "rationale": "<string>"}
}
\end{verbatim}

\section{Rubric --- Five Dimensions}
\label{app:judge-rubric}

All five dimensions share the same 1--5 Likert anchor: $1 = $ very poor, $2 = $ poor, $3 = $ adequate, $4 = $ good, $5 = $ excellent. The fourth dimension, \texttt{hallucination\_inverse}, is inverted internally so that for all five dimensions higher equals better. For each dimension the judge returns both an integer score and a short English rationale (maximum two sentences).

\begin{table}[htbp]
\centering
\small
\begin{tabularx}{\textwidth}{lX}
\toprule
\textbf{Dimension} & \textbf{What is assessed} \\
\midrule
Completeness & Coverage of Art.~30(1) GDPR required elements: controller identity and contact, DPO, purposes of processing, legal basis, categories of data subjects and personal data, categories of recipients, retention periods, and technical/organisational measures (TOMs). Missing elements lower the score. \\
\addlinespace
Scenario Faithfulness & Accuracy with respect to the concrete facts of this scenario: APOR GmbH, Ulm (Mainstraße~67), $\sim$50 employees, internally hosted Microsoft Office and internal servers (no external cloud), no medical/sickness data, 3-month deletion after employee departure, parking authorisation using a distinguishing feature such as a licence plate. Generic ROPA boilerplate lowers the score. \\
\addlinespace
Legal Correctness & Correct use of GDPR terminology and correct citation of legal basis (typically Art.~6(1)(b) for employment data, Art.~6(1)(c) where applicable, Art.~6(1)(f) for parking authorisation). Wrong articles, invented legal bases, or confused controller/processor/DPO roles lower the score. \\
\addlinespace
Hallucination (inverse) & Inverse score for invented facts not supported by the scenario (e.g. invented third-party processors, invented categories such as medical data, invented addresses, retention periods other than 3~months). 5~=~no hallucinations, 1~=~many invented facts. \\
\addlinespace
Overall & Holistic judgement of how useful the document would be as a real Art.~30 ROPA record for this scenario, taking the four dimensions above into account. \\
\bottomrule
\end{tabularx}
\caption{The five judge dimensions. All share the 1--5 Likert anchor; higher is better for every dimension after the hallucination inversion.}
\label{tab:judge-rubric}
\end{table}

\section{Model Configuration and Limitations}
\label{app:judge-config}

The configuration in Table~\ref{tab:judge-config} is taken from the \texttt{judge} block of \texttt{python/judge/prompt.\-json} and the corresponding cells of \texttt{ROPAgen\_JUDGE.ipynb}. It is the minimal information required to reproduce the evaluation layer.

\begin{table}[htbp]
\centering
\small
\begin{tabularx}{\textwidth}{lX}
\toprule
\textbf{Setting} & \textbf{Value} \\
\midrule
Model & \texttt{google/gemini-3.1-pro-preview} (via OpenRouter) \\
Temperature & \texttt{0.0} (deterministic decoding) \\
Output format & JSON object, enforced via OpenRouter \texttt{response\_format = json\_schema} with \texttt{strict: true} \\
Schema name & \texttt{ropa\_judge\_result} \\
Required keys & \texttt{completeness}, \texttt{scenario\_faithfulness}, \texttt{legal\_correctness}, \texttt{hallucination\_inverse}, \texttt{overall} \\
Per-key shape & \texttt{\{"score": int in [1,5], "rationale": string\}} \\
Additional properties & Disallowed at every object level \\
Documents evaluated & 113 ROPA documents (Form~37, Ask~37, Chat~39) \\
Caching & Per-document JSON result cached on disk; reused on re-runs \\
Retry behaviour & Schema-validation or transient API failures retried; persistent failures logged and skipped \\
Logging & Each call logs model, latency, prompt and completion token counts, and raw response payload \\
\bottomrule
\end{tabularx}
\caption{LLM-as-a-judge model configuration.}
\label{tab:judge-config}
\end{table}

\paragraph{Known limitations.} The following limitations of the judge layer are stated for full transparency and discussed further in the Discussion chapter:

\begin{itemize}
\item \textbf{Single-judge bias.} Only one judge model (Gemini~3.1 Pro Preview) operationalises the rubric; inter-rater agreement with a human GDPR expert is not established.
\item \textbf{Model-family bias.} Documents are produced by Mistral Large~3 inside ropagen, while the judge belongs to a different model family (Google). This reduces but does not eliminate stylistic-preference bias.
\item \textbf{Single-scenario bias.} All scores are conditional on the one APOR GmbH scenario; generalisation to other ROPA contexts is not claimed.
\item \textbf{Reference-free.} Unlike the NLP metric layer, the judge does not compare against an expert reference document; it scores against the scenario and Art.~30 directly.
\end{itemize}
